\chapter{Distance}
\label{chap:distance}

% Closed question: How does a finite count of distinctions become an admissible
% notion of distance?

\begin{chapteressay}

A laser rangefinder sends a pulse of light toward a wall. The pulse leaves the
instrument, crosses the room, reflects, and returns. The device measures the
time between departure and return. Multiply half that time by the speed of
light, and the wall has a distance.

No ruler was laid against the wall.

That is the first surprise of distance in physical measurement. Distance looks
like the simplest thing in the world: put a ruler down and count marks. But the
ruler is already an instrument. Its marks are calibrated. Its endpoints must be
aligned. Its temperature matters. Its material expands. The rangefinder makes
the hidden structure explicit. Distance is a count of distinguishable
refinements: light-time ticks, wavelength fringes, ruler marks, odometer pulses,
or some other admissible unit.

The second surprise is that finer distance is not free. If a measurement wants
to distinguish smaller intervals, it must use a finer reference: shorter light
pulses, shorter wavelengths, more stable clocks, more teeth on a wheel, better
phase resolution, more careful calibration. At some point the act of gaining
finer distinction disturbs the system being measured. Distance is therefore not
a passive property waiting untouched. It is a relation between a record, a
reference, and a resolution.

The same structure appears in tailoring. A tailor wraps a tape around a
customer's chest. The tape is flexible, marked in calibrated units, and pulled
with professional judgment: tight enough to remove slack, not so tight that it
compresses the body. The chest measurement is not the body itself. It is the
number of tape intervals admitted under a posture, a tension, and a reference
unit. If the customer inhales, the number changes. If the tape stretches, the
number changes. If the tailor measures at a different height, the number
changes. Distance depends on the admissible comparison.

This chapter names distance as finite refinement count. A measurement record
does not begin with a continuum and then sample it imperfectly. It begins with
distinguishable marks. Distance is what appears when those marks can be ordered,
counted, scaled, and refined. The continuum may appear later as a limit, but the
instrument begins with finite distinctions.

The closed question is:

\begin{center}
\emph{How does a finite count of distinctions become an admissible notion of
distance?}
\end{center}

The answer is that distance is a calibrated count of refinements between
distinguishable events. A laser rangefinder counts light-time. An interferometer
counts fringes. A quantum measurement pays a disturbance cost for finer
position distinctions. A polarizer turns angular separation into transmitted
intensity. An odometer accumulates distance as a sum over speedometer intervals.

Chapter 5 therefore converts the ledger into geometry. Not full geometry yet:
that will require transport, metric structure, and curvature. But the first
physical distance is already enough to change the book. Once marks can be
separated by admissible counts, the ledger has size.

\end{chapteressay}

% Phenomenon arcs use: 1/3 setup -> phenomenon box -> 2/3 structural interpretation.

\section{Distance as Refinement Count}
\label{se:distance-as-refinement-count}

A laser rangefinder measures distance by counting time refinements. If the
round-trip light time is $\Delta t$, the one-way distance is

\[
d = \frac{c \Delta t}{2}.
\]

The factor of two is not cosmetic. The instrument measures a round trip. The
distance to the wall is inferred by dividing the carried record by the geometry
of the experiment. The count is time; the interpretation is distance.

If the instrument resolves one nanosecond, light travels about 30 centimeters
in that time, so the one-way range resolution is about 15 centimeters before
other refinements and corrections. A better clock gives a finer distance. A
shorter pulse can reduce ambiguity. A better detector can identify the returning
edge more accurately. In every case, the distance improves because the
instrument carries finer distinctions.

This means distance is not first a line and only later a number. For the
rangefinder, distance is a count of admissible time intervals under a
calibration of light speed and apparatus geometry. The line is reconstructed
from the count.

The speedometer's odometer performs the same conversion more slowly. Each wheel
rotation carries a calibrated circumference. Each pulse is a fraction of that
circumference. The odometer accumulates those fractions. Distance is not laid
out on the road by the car. It is summed from carried refinements.

This is why distance belongs after comparison. A unit mark alone is not a
distance. Two marks that can be compared and counted under a scale begin to
produce distance. The rangefinder compares send and receive times. The odometer
compares pulses against a circumference. The ruler compares endpoints against
marked intervals. A distance is an ordered count with a named reference.

\section{Interferometry}
\label{se:interferometry}

Interferometry shows how far refinement can go when distance is turned into
phase.

A Michelson interferometer is conceptually simple: a beam splitter, two
mirrors, and a detector. A monochromatic laser beam strikes the beam
splitter, which is a half-silvered glass plate at a 45-degree angle. Half
the light passes through to a mirror at the end of one arm; half is
reflected to a mirror at the end of the other, perpendicular, arm. Both
mirrors reflect their beams back to the beam splitter, where the two
returning beams recombine. The recombined light is then directed to a
photodetector.

The interference at the detector depends on the relative phase of the two
recombined beams. If the two arms have exactly equal optical path length,
the beams add constructively and the detector sees a bright spot. If one
arm is longer by half a wavelength, the beams add destructively (the
phase difference is $\pi$) and the detector sees a dark spot. As the
relative path length is varied continuously, the detector sees alternating
bright and dark fringes.

For light of wavelength $\lambda$, a change of one fringe (bright-dark-bright)
corresponds to a path-length change of $\lambda/2$ in a simple round-trip
geometry. With a $633\,\mathrm{nm}$ helium-neon laser, one fringe
corresponds to about $316.5\,\mathrm{nm}$ of path difference. Distance has
become a count of optical phase distinctions. Counting fringes as one
mirror is moved gives the distance moved in units of $\lambda/2$. A
sub-fringe interpolation can extract a fraction of a fringe by reading
the instantaneous intensity at the detector, giving distance measurements
to nanometer precision.

The metrological significance of interferometry is the chapter's
load-bearing example. The SI unit of length, the meter, was originally
defined as one ten-millionth of the distance from the equator to the
North Pole through Paris (1793), and then as the distance between two
marks on a platinum-iridium bar in S\`evres (1889). In 1960, the
definition was changed to a specific number of wavelengths of a
krypton-86 spectral line, replacing the physical artifact with an atomic
reference. In 1983, the definition was changed again to its current form:
the meter is the distance traveled by light in vacuum in
$1/299{,}792{,}458$ second.

This 1983 definition makes the meter inherently interferometric. To
realize the meter, a metrology institute uses a frequency-stabilized
laser whose wavelength is known to high precision (typically via direct
comparison to atomic transitions). An interferometer then translates
the wavelength into a counted distance. NIST and other national
metrology institutes maintain such instruments; the SI meter is
realized to about one part in $10^{11}$ through this procedure.

The chapter's distance-from-counting structure is precise here. The counted
quantity is fringes; the unit of count is $\lambda/2$; the calibration
certificate is the laser frequency and the vacuum propagation conditions.
The continuous quantity (distance) is the limit of the discrete count as the
fringe spacing becomes infinitesimal relative to the distance being measured.

Interferometry also shows why distance is relational. The instrument measures
path difference. If both arms lengthen by the same amount, the fringe pattern
may not change. If one arm changes relative to the other, the pattern shifts.
The reading is not an isolated length floating in space; it is a comparison
between paths.

The same structure becomes dramatic in gravitational-wave detectors. LIGO is an
interferometer with two long arms. A passing gravitational wave changes the
relative arm lengths by an astonishingly small amount. The detector does not
touch spacetime curvature as an object. It compares optical paths and carries
the differential strain as a fringe-like signal.

The lesson for this chapter is narrower but decisive. Distance becomes
measurable when a finite instrument can refine distinctions. Interferometry
pushes the refinement into phase, where tiny changes produce countable
intensity changes. The result feels continuous because the apparatus is
excellent. Underneath, the record remains a comparison of distinguishable
states.

\section{The Heisenberg Trade-off}
\label{se:the-heisenberg-trade-off}

Finer distinction can cost more than better engineering. In quantum measurement,
the act of measuring position can disturb momentum.

To localize an electron within a distance $\Delta x$, the measurement must use
a probe capable of resolving that scale. A photon with wavelength larger than
$\Delta x$ cannot distinguish features smaller than its own wavelength. A
shorter wavelength means higher momentum. The probe that makes position finer
therefore carries a stronger disturbance.

This is not an afterthought added to quantum theory. It is the distance problem
made sharp: to distinguish smaller intervals, the instrument must bring a finer
reference, and the finer reference is physical. It interacts.

\begin{phenom}{Heisenberg Effect}
\label{ph:heisenberg-effect}

\PhStatement Position refinement and momentum disturbance are coupled. A
measurement that distinguishes a smaller spatial interval requires a probe that
carries enough physical action to disturb the system's momentum record.

\PhOrigin Heisenberg's uncertainty principle identified a structural trade-off
between position and momentum, later formalized as an operator relation and
also visible in operational measurement arguments.

\PhObservation A photon used to locate an electron must have wavelength small
enough to resolve the desired interval. Shorter wavelength means larger photon
momentum, so the measurement that improves $\Delta x$ worsens the momentum
disturbance.

\PhConstraint The instrument may not demand arbitrarily fine distance
distinctions while pretending the reference probe is physically inert.

\PhConsequence Distance is not free resolution. Finer marks require stronger
measurement structure, and that structure can change the record being measured.

\end{phenom}

This phenomenon keeps the distance chapter honest. The rangefinder and
interferometer may make refinement look like a technical challenge only: better
clocks, better lasers, better phase detection. Quantum measurement says that
refinement can become a physical trade-off. The reference used to distinguish
the interval is part of the experiment.

For the speedometer, the analogous cost is mundane: higher tooth counts, faster
electronics, and more aggressive filtering may improve resolution but create
new sensitivity to vibration, slip, or noise. The Heisenberg Effect is the deep
version of a general rule: distance comes from distinction, and distinction
requires an instrument that acts.

\section{Malus Effect}
\label{se:malus-effect}

Distance need not be linear. Sometimes the relevant distance is angular, and
the instrument converts that angular separation into an intensity record.

Polarized light passing through a polarizer gives the clean example. If the
light is linearly polarized at an angle $\theta$ relative to the polarizer's
axis, the transmitted intensity follows

\[
I = I_0 \cos^2 \theta.
\]

At $\theta = 0$, the aligned component passes. At $\theta = 90^\circ$, no ideal
aligned component passes. The polarizer is an angular-distance instrument. It
does not report the angle by drawing it. It reports transmitted intensity, which
is a calibrated function of angular separation.

\begin{phenom}{Malus Effect}
\label{ph:malus-effect}

\PhStatement A physical instrument can measure distance in a space of states by
turning separation from a reference into an admissible intensity record.

\PhOrigin Etienne-Louis Malus discovered that polarized light transmitted
through an analyzer varies with the square of the cosine of the angle between
polarization and analyzer axes.

\PhObservation A polarizer admits the component of the electric field aligned
with its axis. The transmitted intensity records the squared projection, so
angular separation becomes measurable as brightness.

\PhConstraint The instrument may compare only against the reference axis it
carries. It does not report an absolute polarization without a named analyzer
orientation.

\PhConsequence Distance is not only length along a ruler. A calibrated
separation from a reference state can become an admissible distance-like
quantity.

\end{phenom}

This belongs in the distance chapter because it broadens the idea of distance
without abandoning measurement. The polarizer has a reference. The incoming
light has a state relative to that reference. The output is a countable
intensity. A separation in state space becomes a physical reading.

Later gauge theories will make this kind of reference dependence central. For
now, Malus gives a modest lesson: distance is not metaphysical spacing. It is an
admissible comparison against a calibrated reference.

\section{Odometer as Riemann Sum; Spline as Reconstructed Speed}
\label{se:odometer-as-spline}

The odometer integrates the speedometer. The spline enters
elsewhere: not in the odometer itself, but in the reconstruction
of a smooth speed curve from the sampled speed values the
speedometer reports.

If the speedometer reports speeds over small time intervals, the odometer
accumulates distance by summing

\[
\Delta d_i \approx v_i \Delta t_i.
\]

Over a trip, the total distance is

\[
d \approx \sum_i v_i \Delta t_i.
\]

This is a Riemann-sum structure in the dashboard. The car does not need to know
calculus in order to perform it physically. It carries finite pulse counts and
adds calibrated increments.

As the pulse intervals shrink, the sum can approach the familiar integral

\[
d = \int v(t)\,dt.
\]

But the integral is the smooth limit of the finite record, not the first object
the instrument possesses. The odometer has increments. The integral is how
those increments behave under refinement.

This is where distance begins to touch analysis. A variable-speed trip cannot
be captured by one speed reading unless the speed is constant. The odometer
therefore accumulates local contributions. If the contributions are fine enough
and calibrated enough, the sum becomes stable under further refinement. That
stability is what makes distance a trustworthy physical quantity.

The word ``spline'' enters here because the speed between pulses is not
directly recorded. The odometer itself is a Riemann sum, not a spline.
But when the analysis wants a smooth speed curve from the discrete
samples --- to recover $v(t)$ between the pulses --- the reconstruction
is a spline. The instrument or analysis supplies a minimal-curvature
continuation between marks. If the continuation is chosen minimally and
refined under better data, both the speed reconstruction and the
derived distance estimate stabilize. The odometer's sum is the physical
shadow of integration; the spline is the mathematical shadow of speed
reconstruction between samples. The two are separate processes that
share the same underlying refinement discipline.

Chapter 5 can now close. Distance is a finite count of distinctions under a
reference, refined until the count is stable enough for the instrument's
purpose. It may be light-time, phase, quantum localization, angular projection,
or accumulated pulse intervals. In every case, distance is earned by the
ledger.

\begin{bridge}

The chapter's question was how a finite count of distinctions becomes an
admissible notion of distance.

The answer is refinement under calibration. The rangefinder counts light-time
intervals. The interferometer counts phase distinctions. The Heisenberg trade-
off names the cost of finer position resolution. The Malus polarizer turns
angular separation into intensity. The odometer accumulates distance as a sum
over finite speedometer intervals.

The next question is forced. Once distance exists, many continuations remain
possible between measured endpoints. Chapter 6 asks what selects one admissible
continuation as the history actually carried forward.

\end{bridge}

\begin{coda}{A Tailor Measuring a Suit with a Calibrated Tape}

The tailor asks the customer to stand naturally: feet level, shoulders relaxed,
arms at the sides. The tape goes around the chest. It is pulled snugly, but not
tight. The tailor reads the mark, writes the number, and moves to the next
measurement.

The tape is not the body. It is a calibrated distance instrument.

The chest measurement depends on posture, breath, tape tension, fabric, and the
tailor's trained placement. If the customer inhales deeply, the number changes.
If the tape stretches, the number changes. If the tape is held at an angle, the
number changes. The reading is not a metaphysical chest circumference. It is a
carried count of tape intervals under a named practice.

The tailor's tape performs magnitude. The number of centimeters is the carried
quantity. It performs scale. Each mark on the tape is a calibrated unit. It
performs load. The tape must be placed under enough tension to remove slack but
not enough to deform the body. It performs refinement. A tailor can read to the
nearest centimeter, half-centimeter, or millimeter depending on the tape and the
purpose.

There is also a trade-off. To measure the chest, the tailor asks for one
posture. To measure sleeve length, another posture may be needed. To measure a
waist for trousers, the customer's stance and breath matter differently. Trying
to measure every distance at once with the same tape tension would distort the
record. The body is not a rigid object waiting passively for all distances to be
read. The act of measuring selects a reference condition.

The tailor's professional judgment is the admissibility gate. Too loose, and
the tape counts air. Too tight, and the tape counts compression. Wrong height,
and the tape counts the wrong contour. The tailor does not escape this by
claiming exact access to the body. The craft is knowing how to make the
measurement admissible.

When the suit is cut, the numbers become fabric. If the measurements were
compatible and calibrated, the jacket hangs correctly. If one measurement was
taken under a different posture, the error appears in the fit. The tape record,
like the odometer record, becomes trustworthy only because finite distinctions
were counted under a disciplined reference.

The tailor is not doing laboratory physics. Yet the distance structure is the
same. A finite count of calibrated marks becomes an admissible length when the
instrument, posture, tension, and reference are named.

The suit fits because distance was not guessed. It was carried.

\end{coda}
